% GENERATED FILE. Edit canonical lessons, then run npm run generate:books.
% canonical-source-hash: fnv1a64:41458abe495adb13
% canonical-lessons: GU-C207-gadi, GU-C207-teksi, GU-C207-draivar, GU-C207-petrol, GU-C207-bhade-levu

\chapter{Taking a taxi}
\label{ch:gu-taking-a-taxi}
\hlchaptermodality{207}

\begin{chapteropening}
\textbf{By the end of this chapter:} \emph{I can say who drove, what we hired and what the car ran out of.}

Complete the last lesson of chapter 207: I can say who drove, what we hired and what the car ran out of.
\end{chapteropening}

\section[\texorpdfstring{\gu{ગાડી} (gāḍī)}{ગાડી (gāḍī)}]{\gu{ગાડી} (gāḍī) — car, cart}

\label{lesson:GU-C207-gadi}



\begin{quote}
Before the new one: say the Gujarati for stopping place, then the Gujarati for speed.
\end{quote}

\subsection*{You'll want to know: \gu{ગાડી}}
\textbf{\gu{ગાડી}} — \emph{gāḍī} — \textquotedblleft{}car, cart\textquotedblright{}.

\begin{quote}
\textbf{\gu{અમારી} \gu{ગાડી} \gu{નવી} \gu{હતી}.} — \emph{amārī gāḍī navī hatī} — \textquotedblleft{}Our car was new.\textquotedblright{}
\end{quote}

\begin{cousinweb}
\textbf{\gu{ગાડી}} (\emph{gāḍī}) is a tadbhava. CDIAL lists it, glossed \textquotedblleft{}carriage\textquotedblright{}, under \emph{*gāḍḍa}, \textquotedblleft{}cart\textquotedblright{}, a reconstructed form (CDIAL 4116).
\end{cousinweb}

\begin{grammarlens}[title={What you've built}]
One more word for everyday life at A2.
\end{grammarlens}

\subsection*{Your turn}
\begin{itemize}
\raggedright
  \item \emph{Say it:} \emph{gāḍī}
  \item \emph{Say it:} \emph{gāḍī}, once more
  \item \emph{Say it:} \emph{amārī gāḍī navī hatī} — Our car was new.
  \item \emph{Recall:} say the Gujarati for stopping place, then the Gujarati for speed, then say \emph{gāḍī} again
\end{itemize}

\begin{tcolorbox}[breakable,colback=teal!4,colframe=teal!35!black,title={Before you move on}]
What does \gu{ગાડી} mean? (\textquotedblleft{}Car, cart\textquotedblright{}.) Say it once more.
\end{tcolorbox}
\section[\texorpdfstring{\gu{ટેક્સી} (ṭeksī)}{ટેક્સી (ṭeksī)}]{\gu{ટેક્સી} (ṭeksī) — taxi}

\label{lesson:GU-C207-teksi}



\begin{quote}
Before the new one: say the Gujarati for speed, then the Gujarati for car.
\end{quote}

\subsection*{You'll want to know: \gu{ટેક્સી}}
\textbf{\gu{ટેક્સી}} — \emph{ṭeksī} — \textquotedblleft{}taxi\textquotedblright{}.

\begin{quote}
\textbf{\gu{અમે} \gu{સ્ટેશનથી} \gu{ટેક્સી} \gu{લીધી}.} — \emph{ame sṭeśanthī ṭeksī līdhī} — \textquotedblleft{}We took a taxi from the station.\textquotedblright{}
\end{quote}

\begin{cousinweb}
\textbf{\gu{ટેક્સી}} (\emph{ṭeksī}) is a loan from English \emph{taxi}.
\end{cousinweb}

\begin{grammarlens}[title={What you've built}]
One more word for everyday life at A2.
\end{grammarlens}

\subsection*{Your turn}
\begin{itemize}
\raggedright
  \item \emph{Say it:} \emph{ṭeksī}
  \item \emph{Say it:} \emph{ṭeksī}, once more
  \item \emph{Say it:} \emph{ame sṭeśanthī ṭeksī līdhī} — We took a taxi from the station.
  \item \emph{Recall:} say the Gujarati for speed, then the Gujarati for car, then say \emph{ṭeksī} again
\end{itemize}

\begin{tcolorbox}[breakable,colback=teal!4,colframe=teal!35!black,title={Before you move on}]
What does \gu{ટેક્સી} mean? (\textquotedblleft{}Taxi\textquotedblright{}.) Say it once more.
\end{tcolorbox}
\section[\texorpdfstring{\gu{ડ્રાઈવર} (ḍraivar)}{ડ્રાઈવર (ḍraivar)}]{\gu{ડ્રાઈવર} (ḍraivar) — driver}

\label{lesson:GU-C207-draivar}



\begin{quote}
Before the new one: say the Gujarati for car, then the Gujarati for taxi.
\end{quote}

\subsection*{You'll want to know: \gu{ડ્રાઈવર}}
\textbf{\gu{ડ્રાઈવર}} — \emph{ḍraivar} — \textquotedblleft{}driver\textquotedblright{}.

\begin{quote}
\textbf{\gu{ડ્રાઈવરે} \gu{ગાડી} \gu{ધીમે} \gu{ચલાવી}.} — \emph{ḍraivare gāḍī dhīme calāvī} — \textquotedblleft{}The driver drove slowly.\textquotedblright{}
\end{quote}

\begin{cousinweb}
\textbf{\gu{ડ્રાઈવર}} (\emph{ḍraivar}) is a loan from English \emph{driver}.
\end{cousinweb}

\begin{grammarlens}[title={What you've built}]
One more word for everyday life at A2.
\end{grammarlens}

\subsection*{Your turn}
\begin{itemize}
\raggedright
  \item \emph{Say it:} \emph{ḍraivar}
  \item \emph{Say it:} \emph{ḍraivar}, once more
  \item \emph{Say it:} \emph{ḍraivare gāḍī dhīme calāvī} — The driver drove slowly.
  \item \emph{Recall:} say the Gujarati for car, then the Gujarati for taxi, then say \emph{ḍraivar} again
\end{itemize}

\begin{tcolorbox}[breakable,colback=teal!4,colframe=teal!35!black,title={Before you move on}]
What does \gu{ડ્રાઈવર} mean? (\textquotedblleft{}Driver\textquotedblright{}.) Say it once more.
\end{tcolorbox}
\section[\texorpdfstring{\gu{પેટ્રોલ} (peṭrol)}{પેટ્રોલ (peṭrol)}]{\gu{પેટ્રોલ} (peṭrol) — petrol}

\label{lesson:GU-C207-petrol}



\begin{quote}
Before the new one: say the Gujarati for taxi, then the Gujarati for driver.
\end{quote}

\subsection*{You'll want to know: \gu{પેટ્રોલ}}
\textbf{\gu{પેટ્રોલ}} — \emph{peṭrol} — \textquotedblleft{}petrol\textquotedblright{}.

\begin{quote}
\textbf{\gu{ગાડીમાં} \gu{પેટ્રોલ} \gu{ખૂટ્યું}.} — \emph{gāḍīmā̃ peṭrol khūṭyũ} — \textquotedblleft{}The car ran out of petrol.\textquotedblright{}
\end{quote}

\begin{cousinweb}
\textbf{\gu{પેટ્રોલ}} (\emph{peṭrol}) is a loan from English \emph{petrol}.
\end{cousinweb}

\begin{grammarlens}[title={What you've built}]
One more word for everyday life at A2.
\end{grammarlens}

\subsection*{Your turn}
\begin{itemize}
\raggedright
  \item \emph{Say it:} \emph{peṭrol}
  \item \emph{Say it:} \emph{peṭrol}, once more
  \item \emph{Say it:} \emph{gāḍīmā̃ peṭrol khūṭyũ} — The car ran out of petrol.
  \item \emph{Recall:} say the Gujarati for taxi, then the Gujarati for driver, then say \emph{peṭrol} again
\end{itemize}

\begin{tcolorbox}[breakable,colback=teal!4,colframe=teal!35!black,title={Before you move on}]
What does \gu{પેટ્રોલ} mean? (\textquotedblleft{}Petrol\textquotedblright{}.) Say it once more.
\end{tcolorbox}
\section[\texorpdfstring{\gu{ભાડે} \gu{લેવું} (bhāḍe levũ)}{ભાડે લેવું (bhāḍe levũ)}]{\gu{ભાડે} \gu{લેવું} (bhāḍe levũ) — to hire, to rent}

\label{lesson:GU-C207-bhade-levu}



\begin{quote}
Before the new one: say the Gujarati for driver, then the Gujarati for petrol.
\end{quote}

\subsection*{You'll want to know: \gu{ભાડે} \gu{લેવું}}
\textbf{\gu{ભાડે} \gu{લેવું}} — \emph{bhāḍe levũ} — \textquotedblleft{}to hire, to rent\textquotedblright{}.

\begin{quote}
\textbf{\gu{અમે} \gu{એક} \gu{ગાડી} \gu{ભાડે} \gu{લીધી}.} — \emph{ame ek gāḍī bhāḍe līdhī} — \textquotedblleft{}We hired a car.\textquotedblright{}
\end{quote}

\textbf{\gu{ભાડે}}, \emph{bhāḍe}, 'on hire', is a form of \textbf{\gu{ભાડું}}, 'rent'; with \textbf{\gu{લેવું}} it means to hire.

\begin{cousinweb}
\textbf{\gu{ભાડે} \gu{લેવું}} (\emph{bhāḍe levũ}) is built from \emph{bhāḍe}, \textquotedblleft{}on hire\textquotedblright{}, and \emph{levũ}, \textquotedblleft{}to take\textquotedblright{}. CDIAL traces the hire word, Prakrit \emph{bhāḍaya}, \textquotedblleft{}hire\textquotedblright{}, to \emph{*bhārta}, \textquotedblleft{}pertaining to servants\textquotedblright{}, a reconstructed form (CDIAL 9468).
\end{cousinweb}

\begin{grammarlens}[title={What you've built}]
One more word for everyday life at A2.
\end{grammarlens}

\subsection*{Your turn}
\begin{itemize}
\raggedright
  \item \emph{Say it:} \emph{bhāḍe levũ}
  \item \emph{Say it:} \emph{bhāḍe levũ}, once more
  \item \emph{Say it:} \emph{ame ek gāḍī bhāḍe līdhī} — We hired a car.
  \item \emph{Recall:} say the Gujarati for driver, then the Gujarati for petrol, then say \emph{bhāḍe levũ} again
\end{itemize}

\begin{tcolorbox}[breakable,colback=teal!4,colframe=teal!35!black,title={Before you move on}]
What does \gu{ભાડે} \gu{લેવું} mean? (\textquotedblleft{}To hire, to rent\textquotedblright{}.) Say it once more.
\end{tcolorbox}
